\section{Padrões de projeto e práticas (backend)}
\subsection{Separação por camadas (MVC + Service + Repository)}
Os microserviços seguem uma organização típica:
\begin{itemize}[nosep]
  \item \textbf{controller/}: endpoints REST (camada HTTP)
  \item \textbf{service/}: regras de negócio e orquestração
  \item \textbf{repository/}: acesso a dados (Spring Data)
  \item \textbf{entity/}: modelos JPA
  \item \textbf{dto/}: contratos (requests/responses)
  \item \textbf{config/}: configuração Spring (security, rabbit, clients, etc.)
  \item \textbf{messaging/}: consumers/publishers e integração com RabbitMQ (quando aplicável)
  \item \textbf{exception/}: exceções e handlers
  \item \textbf{audit/}: contexto de actor + Envers (em serviços que auditam)
\end{itemize}

\subsection{Proxy Pattern (HTTP declarativo entre microserviços)}
O backend usa clientes HTTP declarativos do Spring (gerados em runtime), seguindo o \textbf{Proxy Pattern}:
\begin{itemize}[nosep]
  \item Interfaces em \texttt{client/NexusClients.java} definem endpoints.
  \item \texttt{config/WebClientConfig.java} constrói \texttt{RestClient} com base URL configurável e gera proxies via \texttt{HttpServiceProxyFactory}.
\end{itemize}
Isto reduz acoplamento e centraliza os contratos de comunicação.

\subsection{Strategy Pattern (ex.: faturação no finance-service)}
O \texttt{finance-service} encapsula o comportamento de emissão de faturas em estratégias:
\begin{itemize}[nosep]
  \item \texttt{InvoiceProviderStrategy} (interface)
  \item Implementações: \texttt{MockInvoiceProviderStrategy}, \texttt{MoloniInvoiceProviderStrategy}, \texttt{NoopInvoiceProviderStrategy}
\end{itemize}
Assim, o provider pode variar por configuração sem alterar a lógica central (\texttt{InvoiceOrchestrator}).

\subsection{Resiliência (Circuit Breaker + Retry)}
O \texttt{sync-service} aplica Resilience4j em chamadas externas (evita cascatas e dá degradação controlada):
\begin{itemize}[nosep]
  \item Circuit Breaker para reduzir falhas repetidas
  \item Retry com backoff para erros transitórios
\end{itemize}

\subsection{Eventos assíncronos e consistência eventual}
\begin{itemize}[nosep]
  \item RabbitMQ para publicação/consumo de eventos (ex.: booking created/status updated).
  \item Fluxos tipo Saga (coreografada) para consistência eventual entre serviços.
\end{itemize}